%==============================================================================
%  (v)  Numerical results and analysis
%
%  First validation, then the analysis proper.  The subsections after
%  Sec. 5.1 are placeholders to be constructed as the results come in --
%  rename and reorder them freely, but keep validation first.
%==============================================================================
\section{Numerical results and analysis}
\label{sec:results}

%==============================================================================
\subsection{Validation}
\label{sec:results:validation}

% Everything here answers one question: is the simulation trustworthy enough
% that the analysis which follows means anything?  Nothing new is claimed.
% Order: global quantities first (one number that can be wrong on its own),
% then profiles, then the quantities the rest of the paper actually relies on.

\subsubsection{Global quantities}
\TODO{The bulk velocity $\Ub/\utau$ against the experimental value, and the
resulting $\Rey$. This is the primary check, because $\Retau$ is imposed and
$\Ub/\utau$ is an outcome: it tests the momentum balance as a whole. Quote the
deviation as a percentage and compare with the spread among the reference
data sets.}

\begin{table}[t]
  \centering
  \caption{Global validation. \TODO{Present $\Ub/\utau$, $\Umax/\Ub$ and the
  secondary-current magnitude for each case, against experiment and against the
  earlier LES.}}
  \label{tab:validation}
  \small
  \begin{tabular}{lccc}
    \toprule
    Quantity & Present & Experiment & Deviation \\
    \midrule
    \TODO{} & & & \\
    \bottomrule
  \end{tabular}
\end{table}

\subsubsection{Mean flow and secondary currents}
\TODO{Isovels in the cross-section with the secondary-current vectors, against
the reference figures. Then the lateral distribution of depth-averaged velocity
and of bed shear stress. Say what is captured (the junction vortex pair, the
velocity dip, the shear-stress maximum) and, equally, what is not.}

\begin{figure}[t]
  \centering
  % \includegraphics[width=\columnwidth]{isovels}
  \framebox[\columnwidth][c]{\rule{0pt}{3.4cm}\TODO{isovels + secondary currents}}
  \caption{\TODO{Mean streamwise velocity and in-plane vectors; present against
  the reference data in the same axes and normalisation.}}
  \label{fig:isovels}
\end{figure}

\subsubsection{Turbulence statistics}
\TODO{Vertical profiles of $\meanx{U}$ in wall units at the main-channel,
junction and floodplain stations, and the Reynolds stresses. State the sign
convention used for $\avg{\fluc{u}\fluc{v}}$ explicitly, since it differs
between the reference works.}

\subsubsection{Adequacy of the sampling for the spectral analysis}
\TODO{The check the later sections depend on, and which is easy to omit: that
the record is long enough for the SPOD estimator. Report the number of
independent blocks, the resulting confidence interval on $\spodev_1$, and the
two-point coherence showing that $L_x$ is long enough that the largest resolved
$\kx$ modes are not constrained by the box.}

%==============================================================================
% The analysis proper.  Placeholders -- rename as the results take shape.
% Suggested spine, one subsection per claim, each claim falsifiable:
%   5.2  what the SPOD spectra show           (which (k_x, omega) carry energy)
%   5.3  what the leading modes look like     (the structures themselves)
%   5.4  what the resolvent predicts          (gains, and the same modes)
%   5.5  where the two agree, and what the disagreement means
%   5.6  dependence on the depth ratio D_r
%==============================================================================

\subsection{\TODO{Energy distribution in the wavenumber--frequency plane}}
\label{sec:results:spectra}

\TODO{To be constructed. The SPOD eigenvalue spectra $\spodev_j(\kx,\omega)$:
where the energy sits, whether the leading mode separates from the rest (the
low-rank question), and at which $(\kx,\omega)$ that separation is largest.}

%------------------------------------------------------------------------------
\subsection{\TODO{Leading SPOD modes}}
\label{sec:results:modes}

\TODO{To be constructed. The structures themselves at the $(\kx,\omega)$
identified above, and their physical identification --- junction shear-layer
rollers, vertical-axis surface vortices, secondary-current modes.}

%------------------------------------------------------------------------------
\subsection{\TODO{Resolvent gains and response modes}}
\label{sec:results:gains}

\TODO{To be constructed. $\gaink_1(\kx,\omega)$ over the same plane as the SPOD
spectra, plotted on the same axes and the same colour scale so that the two can
be read against each other.}

%------------------------------------------------------------------------------
\subsection{\TODO{Comparison of SPOD and resolvent modes}}
\label{sec:results:comparison}

\TODO{To be constructed. The alignment $\abs{\ip{\spodmode_1}{\resu_1}}$ as a
function of $(\kx,\omega)$. Frame the disagreement as a measurement of the
forcing colour, not as a failure.}

%------------------------------------------------------------------------------
\subsection{\TODO{Effect of the depth ratio}}
\label{sec:results:dr}

\TODO{To be constructed. How the above changes with $\Dr$, and whether any
scaling collapses the cases.}
